\section{Related work}
\label{sec:related}

The relevant comparisons concern three different tasks: visiting every vertex or
edge, meeting at one vertex, and the distance-one contact required here. They also
involve different ways of supplying direction information and different orders of
quantification over controllers, graphs, labellings and starts. We organize the
comparison around these distinctions, rather than treating all finite-state walkers
as one model.

\paragraph{Labyrinth automata, quantitative traps, and universal hosts.}
Budach's work on labyrinth exploration~\cite{Budach1978} and the universal-trap
paper of Antelmann, Budach and Rollik~\cite{AntelmannBudachRollik1979} are central
historical predecessors. For their historical theorem statements we rely on the
account of Kilibarda, Kudryavtsev and U\v{s}\'{c}umli\'{c}~\cite{KilibardaKudryavtsevUscumlic2003},
not on uninspected originals. Hemmerling's monograph explicitly treats traps with a
specified starting position and estimates their size in terms of internal
states~\cite[pp.~133--134]{Hemmerling1989}. Thus neither state-dependent trap size
nor a controller-dependent adversary is a new organizing principle.

The observation interface can also be quite close. Hoffmann's inspected opening
pages define plane mazes as full sublabyrinths of the square lattice, with free
compass directions observed locally, and state a one-pebble exploration
impossibility~\cite[pp.~433--434]{Hoffmann1981}. His finite-automaton definition
uses a legal move at every step. Induced geometry and compulsory motion therefore
do not, by themselves, separate our problem from that literature. What must still
be transferred is the objective: an unvisited vertex is not a certificate that two
moving copies always remain nonadjacent.

A quantitative predecessor makes this limitation especially concrete.
The survey~\cite[Theorem~4, p.~25]{KilibardaKudryavtsevUscumlic2003} attributes to
Antelmann a proper planar exploration trap of size at most
$C\exp\sqrt{2s\log_2(2s)}$ for an initialized $s$-state automaton. The cited statement
does not specify a numerical $C$. Antelmann's original article and dissertation
remain uninspected; neither an exact low-state constant nor a rendezvous conclusion
is inferred from this asymptotic bound.

Universal traps require an equally explicit qualification. Kilibarda proves
universal exploration traps for finite families of automata~\cite[Theorems~1.2--1.3]{Kilibarda1990}.
His later reduction paper gives a finite mosaic trap for each finite collective of independent automata and,
in a separate theorem, one \emph{infinite} mosaic trap for all such
collectives~\cite[Theorems~3.5 and~3.8]{KilibardaReduction2021}. Consequently our
finite obstruction menu and fixed host are not presented as the invention of
universal trapping. The specific assembly preserves the complete masks queried by
two executions, produces an explicitly bounded finite induced subcubic grid tree,
and permits only the bad starts to depend on the controller. Its failure predicate
is perpetual nonadjacency. The earlier exploration theorems are genuine quantifier
predecessors, but none of the cited statements supplies that entire conclusion.

\paragraph{Finite-state exploration, exact small memory, and periodic traversal.}
Graph exploration supplies both asymptotic and exact-state comparisons. The work
of Fraigniaud, Ilcinkas, Peer, Pelc and Peleg~\cite{FraigniaudEtAl2005} studies
port-based automata; a modern account of the space-complexity setting is given by
Disser and Klimm~\cite{DisserKlimm2022}. The latter also makes explicit why vertices
with the same observed incident colours are indistinguishable. This is a useful
neighbor of our repeated-mask requirement, although exploration is a one-agent
objective and its label conventions are not an absolute compass.

The closest numerical trap bound is particularly important to acknowledge.
Ilcinkas~\cite[Theorem~2.1, p.~52]{Ilcinkas2006} constructs, for every $K$-state
port-based Moore automaton and $d\ge3$, a simple-graph exploration trap with at most
$K+d+2$ vertices and maximum degree $d$. Substituting $K=2,d=3$ gives seven.
The construction leaves a \emph{vertex} unvisited, so an edge-versus-vertex slogan
would not distinguish it from our theorem. Instead, a transfer must preserve the
controller interface and induced-grid geometry and must create a nonadjacent pair
of executions. No such state-preserving transfer is established here.

Exact one- and two-state classifications are also already available, rather than
being a new kind of question. Fraigniaud, Ilcinkas and
Pelc~\cite{FraigniaudIlcinkasPelc2008} classify regular graphs explorable by
port-Mealy automata: the one-state class consists of the one-vertex graph, $K_2$ and
cycles, while the two-state class consists of cliques and cycles. An automaton for
a graph must work for every port assignment and initial position. These are strong
predecessors for the small-memory viewpoint, but not determinations of our
max--min trap parameters for a fixed compass-mask interface.

For the structural results, Diks, Fraigniaud, Kranakis and
Pelc~\cite{DiksEtAl2004} already use a degree-based state lower bound and the
cyclic next-port rule for perpetual tree exploration. Kilibarda's
Proposition~1.1~\cite{Kilibarda1990} states that every rotation system of a tree
induces one contour; the same paper gives exact small-state exploration thresholds
in a richer local-vision model. Thus the contour mechanism in
Section~\ref{sec:paths} is classical. Our four-state table stores the incoming
\emph{compass} port internally, because it is not supplied as an observation;
the radius-one rendezvous conclusion additionally uses path phase contact.
This is not a per-vertex rotor-memory model: the fixed tree has no writable local
state.

Periodic exploration is another close, but different, extremal problem.
Ilcinkas~\cite{Ilcinkas2008Ports} gives a three-state explorer with period at most
$4n-2$ after a suitable port assignment. That problem minimizes a guaranteed
exploration period while permitting the designer to choose local orientations;
its configurations include an incoming port. Our sharp-period theorem instead
maximizes the least tagged period over allowed compass paths and controllers,
under an internal-state budget. Neither the periodic-exploration bound nor the
standard contour alone gives the weighted capacity and sharp all-state formula
proved in Section~\ref{sec:support-structure}.

\paragraph{Meeting identical agents: resources and port quantifiers.}
Pelc's survey provides the broader deterministic network-rendezvous
context~\cite{Pelc2012Survey}. A particularly relevant compass-labyrinth result is
Kilibarda's type-meeting theorem~\cite[Theorem~10]{Kilibarda2021}. Its strong version
uses two copies of one collective and same-vertex contact; the minimal universal
types for finite mosaic labyrinths are $(1,2)$ and $(2,0)$. A type counts active
automata and pebbles, not internal states. In particular, copying a collective of
type $(2,0)$ is not a theorem about two single agents. The negative side concerns
same-vertex meeting and does not supply a distance-one trap or the exact values
four, seven or eight in our model.

Bounded-memory rendezvous in port-labelled graphs is more developed than a
comparison based only on a few constants would suggest. Czyzowicz, Kosowski and
Pelc~\cite{CzyzowiczEtAl2012} establish log-space rendezvous bounds in arbitrary
graphs. For trees, the distinction between the following two domains is essential.
Fraigniaud and Pelc~\cite[Definitions~1.1--1.2 and Section~1.2]{FraigniaudPelc2011}
require success for every port labelling and restrict admissible initial pairs by
perfect symmetrizability. Their simultaneous-start result does not cover every pair
that happens to be nonsymmetric under one fixed labelling. Czyzowicz, Kosowski and
Pelc~\cite{CzyzowiczEtAl2014} prove a logarithmic memory lower bound even for
simultaneous starts on lines in that latter, broader domain. The inspected author
versions explicitly expose the correction to the earlier conference formulation.
These are not conflicting bounds once the domains are kept distinct.

Those models supply arrival-port information, permit waiting, and require
same-node meeting. Our compass, compulsory motion, identical initial states and
integer-time distance-one event must therefore all be retained when comparing the
path threshold or the still-open three-state case. In particular, a lower bound on
same-node rendezvous is not automatically a lower bound on distance-one contact,
even if it already applies to simultaneous starts.

\paragraph{Recent nearby models.}
The 2026 preprint of Das and Pelc~\cite{DasPelc2026} is a close current finite-state
comparison: it asks for a universal deterministic automaton on all feasible
rendezvous instances of port-labelled graphs. Its negative result with finitely
many stationary pebbles already has simultaneous-start line witnesses.
Thus start delay alone cannot explain the separation from our positive path
result. The same-node predicate, port observations and feasibility quantifiers
remain decisive; the preprint does not settle our three-state question.

Recent grid and neighborhood work illustrates two other boundaries.
Gao and Pelc~\cite{GaoPelc2025} study identical agents with a coherent compass on
the infinite grid, but allow waiting and permanent footprints, and require a common
node. Eguchi, Kitamura and Izumi's \emph{neighborhood rendezvous}~\cite{EguchiEtAl2022}
assumes initially adjacent agents and still seeks same-node meeting; it does not
mean success at distance one. Its randomized, labelled-graph setting is also not a
fixed small-state anonymous controller. These papers clarify model choices rather
than provide numerical predecessors to our trap thresholds.

\paragraph{Graph-walking realization and deterministic synthesis.}
Local descriptions of finite-state computations have strong precedents.
Martynova~\cite{Martynova2022} encodes graph-walking computations by attaching
incoming and outgoing state sets to directions and imposing node consistency.
Theorem~5 of the inspected author version transforms a \emph{given} automaton
into a signature describing accepted graphs. Here a graph and a cyclic scenario
are prescribed, while the common transition table is unknown. This reverses the
synthesis direction; it does not make the use of local state annotations itself
new.

Similarly, Ulyantsev, Buzhinsky and Shalyto~\cite{UlyantsevBuzhinskyShalyto2016}
identify deterministic machines from scenarios using state assignments and
transition consistency. Our occurrence-colouring formulation belongs naturally to
this exact-synthesis neighborhood. The content specific to
Theorem~\ref{thm:three-state-realization} is a normal form for a restricted recurrent
scenario: arbitrary row outputs are eliminated in favor of a finite common
mask/difference profile and one bit per internal vertex. It is not a general theorem
that scenario synthesis is new or that local annotations guarantee global
realizability without further constraints.

\paragraph{Algebraic consistency and the scope of the contribution.}
Signed-graph balance and switching are classical; see
Zaslavsky~\cite[Section~2, Proposition~2.1]{Zaslavsky1982}. Gain graphs generalize
edge signs to group elements, with identity product around balanced
cycles~\cite{Zaslavsky1989}. This is the correct context for our signed-graph
special cases, but not for the entire realization criterion. A profile system may
have rows involving more than two bits; it balances all linear dependencies of one
matrix. Successor differences may be zero, and the existential choice of one common
profile can produce non-affine phase projections. Replacing the full criterion by
one signed graph or by permutation-only transport would change the theorem.

Nor is polynomial recognition at a fixed state budget and alphabet asserted as a
new complexity result: the set of complete controller tables is already finite
independently of path length. The contribution claimed here is the proved
combination of the exact model-specific trap thresholds, weighted recurrent-support
bounds with sharp attainment, the special affine realization normal form, and
observation-preserving obstruction assembly. The inspected predecessors do not
supply this combination under the same hypotheses. That comparison is not an
absolute priority claim; the access limitations on historical originals remain
explicit.
